%----------------------------------------------------------------------------------------
% Tailored CV — Doktorand:in, Mehrdimensionale OMICS-Datenanalyse, ISAS Dortmund
%----------------------------------------------------------------------------------------

\documentclass[10pt,a4paper,sans]{moderncv}

\usepackage[utf8]{inputenc}
\usepackage[T1]{fontenc}

\moderncvstyle{classic}
\moderncvcolor{blue}

\usepackage[scale=0.78]{geometry}
\usepackage{ragged2e}

%----------------------------------------------------------------------------------------

\firstname{Kundai Farai}
\familyname{Sachikonye}

\address{Auf der Lichtung 19}{82178 Puchheim, Germany}
\mobile{(+49) 1626386344}
\email{kundai.sachikonye@bitspark.com}
\homepage{github.com/fullscreen-triangle}{github.com/fullscreen-triangle}

%----------------------------------------------------------------------------------------

\begin{document}

\makecvtitle

%----------------------------------------------------------------------------------------
%	EDUCATION
%----------------------------------------------------------------------------------------

\section{Education}

\cventry{2019--2021}{PhD candidate, Computational Lipidomics}{Technische Universität München}{}{}{
  Mass spectrometry pipelines, multi-omics statistical modelling, federated learning
  for distributed clinical analysis. Departed to pursue independent computational research.
}
\cventry{2018--2019}{MSc Computational Biology}{Universität Tübingen}{}{}{}
\cventry{2014--2017}{MSc Pharmaceutical Biotechnology}{HAW Hamburg}{}{}{}
\cventry{2011--2014}{BSc Biochemistry and Cell Biology}{Jacobs University Bremen}{}{}{}

%----------------------------------------------------------------------------------------
%	RESEARCH OUTPUT
%----------------------------------------------------------------------------------------

\section{Research Output}

\cvitem{2025}{
  \textbf{On the Thermodynamic Consequences of Categorical Completion on Biological
  Systems.}
  MHC binding affinity correlates with categorical richness ($r=0.73$,
  $p<10^{-12}$, $n=2{,}847$ IEDB peptides). AUC 0.81 vs NetMHCpan 0.68.
  \textit{Manuscript.}
}

\cvitem{2025}{
  \textbf{On Disease Epistemology in Blind Receiver.}
  Loop-holonomy self-consistency diagnostic; AUC$=1.00$; 25/25 validation experiments;
  sparse $\ell_1$ LP for personalised therapeutic target selection.
  \textit{Manuscript. AIMe Registry / TUM Weihenstephan.}
}

\cvitem{2025}{
  \textbf{Syndrome: Categorical Resolution of Disease Trajectories.}
  Eight oscillator classes; disease vector $\mathbf{D}\in[0,1]^8$;
  $\mathcal{O}(N\log N)$ trajectory computation; personalised intervention via
  sparse LP. \textit{Manuscript.}
}

%----------------------------------------------------------------------------------------
%	RESEARCH SOFTWARE — OMICS AND IMAGING
%----------------------------------------------------------------------------------------

\section{OMICS and Imaging Tools}

\cvitem{gospel}{
  \textbf{Multi-Omics Genomic Analysis.}
  $10^6$ variants/second VCF processing; RNA-seq integration; Bayesian networks,
  GMMs, variational Bayes. Validated: 1000 Genomes Ph.\ 3, gnomAD v3.1.2, UK Biobank.
  Pathway analysis (KEGG, Reactome), protein structure (RDKit). Live:
  \href{https://vivid-symbolism.vercel.app/search}{vivid-symbolism.vercel.app/search}.
  \href{https://github.com/fullscreen-triangle/gospel}{github.com/fullscreen-triangle/gospel}
}

\cvitem{lavoisier}{
  \textbf{MALDI MSI and LC-MS/MS Analysis.}
  Dual pipeline: numerical MS1/MS2 + CNN visual pipeline (PyTorch);
  98.9\,\% NIST accuracy; Ray distributed; $>$100\,GB mzML support.
  Virtual instrument: \href{https://lavoisier.vercel.app/experiment}{lavoisier.vercel.app/experiment};
  Rust executable for local large-scale runs.
  \href{https://github.com/fullscreen-triangle/lavoisier}{github.com/fullscreen-triangle/lavoisier}
}

\cvitem{helicopter}{
  \textbf{Multi-Modal Microscopy Analysis.}
  Sequential constraint exclusion across fluorescence, brightfield, phase-contrast,
  and spectral channels; twelve modalities; effective sub-nanometre resolution.
  Live (GPU, browser, TIFF):
  \href{https://hieronymus-omega.vercel.app/observe/microscopy}{hieronymus-omega.vercel.app/observe/microscopy}.
  \href{https://github.com/fullscreen-triangle/helicopter}{github.com/fullscreen-triangle/helicopter}
}

\cvitem{syndrome}{
  \textbf{Personalised Disease Trajectory and Risk Assessment.}
  Disease vector $\mathbf{D}\in[0,1]^8$ across eight cellular oscillator classes;
  $\mathcal{O}(N\log N)$ trajectory to healthy attractor; sparse $\ell_1$ LP for
  minimum-side-effect therapeutic target selection. Validated against established
  biomarkers.
  \href{https://github.com/fullscreen-triangle/syndrome}{github.com/fullscreen-triangle/syndrome}
}

\cvitem{shakespear}{
  \textbf{Protein Analysis Instrument.}
  Folding, trajectory, catalysis, and dynamics via Kuramoto oscillator synchronization
  in S-entropy space. Applicable to proteomics readouts and protein-state transitions
  in multi-omics pipelines. Browser-based, all computation client-side. Live:
  \href{https://shakespear-nine.vercel.app/instrument}{shakespear-nine.vercel.app/instrument}.
  \href{https://github.com/fullscreen-triangle/shakespear}{github.com/fullscreen-triangle/shakespear}
}

\cvitem{honjo-masamune}{
  \textbf{Categorical Cheminformatics.}
  GPU-accelerated zero-parameter molecular spectroscopy: vibrational modes, docking,
  3D structure, fingerprints, molecular network analysis. 39 NIST compounds; S-entropy
  coordinate framework; all computation browser-side. Live:
  \href{https://honjo-masamune.vercel.app/tools}{honjo-masamune.vercel.app/tools}.
}

%----------------------------------------------------------------------------------------
%	RESEARCH SOFTWARE — AI AND LLM TOOLING
%----------------------------------------------------------------------------------------

\section{AI and LLM Tools}

\cvitem{purpose}{
  \textbf{Domain-Specific Model Training --- Backward Training Principle.}
  Formal proof that training at the pure-intent regime (hardest single-objective
  instance) yields competence on all constrained sub-regimes by inclusion.
  $(N{+}1)\times$ to $2^{N-1}\times$ sample savings vs.\ conventional curriculum.
  Aperture Foundation Model + Domain Connector + Resolver Registry.
  \href{https://github.com/fullscreen-triangle/purpose}{github.com/fullscreen-triangle/purpose}
}

\cvitem{four-sided-triangle}{
  \textbf{Multi-Model RAG and Domain Expert Extraction.}
  8-stage metacognitive pipeline; adaptive hybrid LLM and mathematical solver selection;
  Rust backend (10--50$\times$ speedup); FastAPI + React; Docker/Kubernetes.
  \href{https://github.com/fullscreen-triangle/four-sided-triangle}{github.com/fullscreen-triangle/four-sided-triangle}
}

%----------------------------------------------------------------------------------------
%	RESEARCH EXPERIENCE
%----------------------------------------------------------------------------------------

\section{Research Experience}

\cventry{2021--present}{Independent AI Systems Developer}{Bitspark GmbH}{}{}{
  Multi-omics analysis frameworks, LLM orchestration, domain-specific model training,
  agentic pipelines, and scientific data tools. All systems publicly deployed.
}

\cventry{2019--2021}{Computational Lipidomics}{TUM / Bayerisches Forschungsinstitut}{Munich}{}{
  Mass spectrometry pipelines (peak detection, ion annotation, feature engineering),
  multi-omics statistical modelling, federated learning for distributed clinical
  analysis. Python, Java.
}

\cventry{2017--2018}{Glycomics and Proteomics}{Universitätsklinikum Hamburg-Eppendorf}{}{}{
  Automated proteomics and glycomics pipeline for LC-MS/MS and MALDI-TOF data.
}

\cventry{2013}{Cancer Biomarker Discovery}{University Medical Center Groningen}{}{}{
  UPLC-MS/MS shotgun proteomics for cancer biomarker discovery; statistical analysis.
}

%----------------------------------------------------------------------------------------
%	TECHNICAL SKILLS
%----------------------------------------------------------------------------------------

\section{Technical Skills}

\cvitem{OMICS / Bioinformatics}{
  Variant calling (VCF), RNA-seq quantification, sequence alignment, single-cell
  analysis, pathway analysis (KEGG, Reactome, GSEA), protein structure (RDKit),
  multi-omics integration (MOFA, UMAP)
}

\cvitem{Imaging}{
  Fluorescence, brightfield, phase-contrast, spectral microscopy; MALDI MSI;
  multi-modal image fusion (OpenCV); image segmentation (PyTorch CNNs)
}

\cvitem{ML / DL}{
  PyTorch, scikit-learn, XGBoost; CNNs, LSTMs, Transformers, GNNs;
  Bayesian networks, Gaussian mixture models, variational Bayes;
  LoRA and parameter-efficient fine-tuning; RAG pipelines
}

\cvitem{LLM / AI}{
  Claude (Anthropic), OpenAI, Ollama, Replicate, Together AI; multi-provider
  orchestration; domain-specific fine-tuning; retrieval-augmented generation;
  tool use and function calling
}

\cvitem{Languages}{Python, R, Rust, TypeScript, Java, C, SQL, Bash}

\cvitem{Infrastructure}{
  Ray, Dask, Docker, FastAPI, AWS, Vercel, PostgreSQL, Git,
  memory-mapped I/O for large-scale datasets
}

\cvitem{Mass Spectrometry}{
  LC-MS/MS, MALDI-TOF / MALDI MSI, UPLC-MS/MS, TIMS-PASEF, QTOF
}

%----------------------------------------------------------------------------------------
%	LANGUAGES
%----------------------------------------------------------------------------------------

\section{Languages}

\cvitemwithcomment{English}{Native / Fluent}{}
\cvitemwithcomment{German}{Conversational}{Living and working in Germany since 2011}

%----------------------------------------------------------------------------------------

\renewcommand{\listitemsymbol}{-~}

\end{document}
